\documentclass[11pt]{article}

\usepackage[a4paper,margin=1in]{geometry}
\usepackage[T1]{fontenc}
\usepackage[utf8]{inputenc}
\usepackage{amsmath,amssymb,amsthm}
\usepackage{mathtools}
\usepackage{enumitem}
\usepackage{booktabs}
\usepackage{xcolor}
\usepackage{algorithm}
\usepackage{algpseudocode}
\usepackage{hyperref}
\usepackage{cleveref}

\hypersetup{colorlinks=true,linkcolor=blue!50!black,citecolor=blue!50!black,urlcolor=blue!50!black}

\theoremstyle{plain}
\newtheorem{theorem}{Theorem}[section]
\newtheorem{lemma}[theorem]{Lemma}
\newtheorem{proposition}[theorem]{Proposition}
\newtheorem{corollary}[theorem]{Corollary}
\theoremstyle{definition}
\newtheorem{definition}[theorem]{Definition}
\newtheorem{example}[theorem]{Example}
\theoremstyle{remark}
\newtheorem{remark}[theorem]{Remark}

\crefname{theorem}{Theorem}{Theorems}
\crefname{lemma}{Lemma}{Lemmas}
\crefname{proposition}{Proposition}{Propositions}
\crefname{corollary}{Corollary}{Corollaries}
\crefname{definition}{Definition}{Definitions}
\crefname{remark}{Remark}{Remarks}

\newcommand{\CP}{\mathbb{CP}}
\newcommand{\R}{\mathbb{R}}
\newcommand{\type}{\mathrm{type}}
\newcommand{\level}{\mathrm{level}}
\newcommand{\msplit}{\mathrm{split}}
\newcommand{\weld}{\mathrm{weld}}
\newcommand{\st}{\mathrm{st}}
\newcommand{\pos}{\mathrm{pos}}
\setlength{\emergencystretch}{2em}

\title{\bf Stellar subdivision schemes that split $2$-faces,\\ and typed-cycle hybrids with Maubach bisection}
\author{Research note (not part of the article)}
\date{30 September 2026}

\begin{document}
\maketitle

\begin{abstract}
We record a stellar subdivision scheme for $4$-dimensional ordered meshes
that splits a $2$-face (a triangle) of each cell instead of an edge
(Maubach) or a facet (Mitchell). With the right rule for the types of the
children, the scheme is combinatorially correct: its invariants are
Maubach's with the split edge $\{0,l\}$ replaced by the split face
$\{0,1,l\}$, and the proof of the article goes through verbatim. The
scheme belongs to a family of \emph{typed cycles}, for which we give three
local conditions that are sufficient, and in exhaustive tests also
necessary, for the invariants to be preserved. Geometrically, however, a
scheme that splits only faces of dimension $\ge2$ never removes an edge,
so cell diameters never tend to zero. Hybrid cycles mixing edge and face
splits do shrink cells, but none of those we tested keeps shape
regularity. All claims marked \emph{verified} were checked by the
scripts in \texttt{scripts/}.
\end{abstract}

\section{Setting}

We use the notation and framework of the article
(\emph{Stellar subdivision schemes and pointerless Maubach bisection on
balanced seeds}, \S3). A cell is an ordered $d$-simplex
$\sigma=\langle\sigma[0],\dots,\sigma[d]\rangle$, $\partial_i$ deletes the
$i$-th vertex, and equalities $\partial_i\sigma=\partial_j\omega$ are
equalities of \emph{ordered} simplices. An ordered stellar subdivision
$(\delta,v)_l$ of a face $\delta=\sigma[p_0,\dots,p_k]$ replaces every cell
$\sigma\supseteq\delta$ by its $k+1$ children
\[
  c_p\sigma \;=\; \text{$\sigma$ with $\sigma[p]$ deleted and $v$ inserted at position $l$},\qquad p\in\{p_0,\dots,p_k\},
\]
so that
\begin{equation}\label{eq:outer}
  \partial_l\,c_p\sigma=\partial_p\sigma .
\end{equation}
A \emph{seed} is \emph{strongly compatible} if all its cells have level $0$
and $\partial_i\sigma=\partial_j\omega$ implies $i=j$. Equivalently, the
seed is balanced and every cell is ordered by color. Examples are Kuhn's
cube and Gaifullin's $\CP^2_{15}$.

\section{The face scheme}\label{sec:face}

Let $d=4$ and $\sigma=\langle a,b,c,d,e\rangle$. The proposal assigns to
each cell an integer $l(\sigma)\in\{4,3,2\}$ (equivalently a type in
$\{0,1,2\}$), splits the $2$-face
\[
  F(\sigma)=\{0,1,l(\sigma)\}
\]
(that is, $\langle a,b,e\rangle$, $\langle a,b,d\rangle$ or
$\langle a,b,c\rangle$), and inserts the new vertex $v$ at position
$l(\sigma)$:
\begin{center}
\begin{tabular}{@{}ccl@{}}
\toprule
$l$ & split face & children \\
\midrule
$4$ & $\langle a,b,e\rangle$ & $\langle a,b,c,d,v\rangle,\ \langle a,c,d,e,v\rangle,\ \langle b,c,d,e,v\rangle$\\
$3$ & $\langle a,b,d\rangle$ & $\langle a,b,c,v,e\rangle,\ \langle a,c,d,v,e\rangle,\ \langle b,c,d,v,e\rangle$\\
$2$ & $\langle a,b,c\rangle$ & $\langle a,b,v,d,e\rangle,\ \langle a,c,v,d,e\rangle,\ \langle b,c,v,d,e\rangle$\\
\bottomrule
\end{tabular}
\end{center}

\paragraph{The type of the children.}
The proposal does not fix $l$ for the children. There are three
level-independent choices, and only one of them works (verified).

\begin{center}
\begin{tabular}{@{}ll@{}}
\toprule
rule for the children & result on the reflected Kuhn torus $\mathbb{Z}_4^4$ ($6144$ cells)\\
\midrule
same $l$ as the parent                   & incompatible after one uniform round ($12288$ of $18432$ cells)\\
$l\mapsto l+1$ (i.e.\ $2\to3\to4\to2$)     & incompatible after one uniform round\\
$l\mapsto l-1$ (i.e.\ $4\to3\to2\to4$)     & compatible (\Cref{thm:face})\\
\bottomrule
\end{tabular}
\end{center}
So, as in Maubach's scheme, $l$ decreases by one per level and wraps
around:
\[
  l(\sigma)=4-(\level(\sigma)\bmod 3).
\]
This is Maubach's rule
$\type(\sigma)=d-(\level(\sigma)\bmod d)$ with the \emph{prefix} $\{0\}$ of
the split edge $\{0,l\}$ replaced by the prefix $\{0,1\}$.

\begin{definition}[face invariants]\label{def:face}
A $4$-mesh satisfies the \emph{face invariants} if it satisfies
(OneLevel) and, for adjacent cells with $\partial_i\sigma=\partial_j\omega$,
\begin{align*}
  \level(\sigma)=\level(\omega)&\;\Longrightarrow\; i=j\ \text{ or }\ \{i,j\}\subseteq F(\sigma), \tag{A}\\
  \level(\sigma)=\level(\omega)+1&\;\Longrightarrow\; i=l(\omega)\ \text{ and }\ j\in F(\omega). \tag{B}
\end{align*}
\end{definition}
In (B), $i=l(\omega)$ is the position of the weld vertex of the finer
cell. A strongly compatible seed satisfies the face invariants.

\begin{theorem}\label{thm:face}
On a conforming $4$-mesh satisfying the face invariants, the Maubach
refinement algorithm (\textsc{MaubachSubdivide} in the article), with the split edge
$\sigma[0,\type(\sigma)]$ replaced by the split face $\sigma[F(\sigma)]$,
the loop over $i\in\{1,\dots,d\}\setminus\{\type(c)\}$ replaced by
$i\notin F(c)$, and the new vertex inserted at $l(\sigma)$, terminates and
only splits cells of level at most $\level(\sigma)$. It produces a
conforming mesh satisfying the face invariants, and it is consistent,
alternating and level-homogeneous.
\end{theorem}

This is the case $d=4$, $k=2$ of \Cref{thm:prefix}, which in turn is a
special case of \Cref{thm:cycle}. The proof is in \Cref{sec:cycles}.

\paragraph{Empirical invariants.}
Under uniform and adaptive refinement from the Kuhn torus (up to level
$10$), the exceptional pairs in (A) are more restricted than (A) allows.
With $l=4$ only $\{i,j\}=\{0,1\}$ occurs, and with $l\in\{3,2\}$ only
$\{0,l\}$ and $\{1,l\}$ occur. Every shared facet is an equality of ordered
simplices. The weaker statement (A) is the one that is preserved
inductively, so it is the one to state.

\paragraph{Why the other rules fail.}
In the language of \Cref{sec:cycles}, the two rejected rules fail for
different reasons.
\begin{itemize}
\item With $l$ fixed, (C1) fails: the weld position $l$ belongs to the
  children's split face $\{0,1,l\}$, so a child would have to split a
  face containing its own weld vertex.
\item With $l\mapsto l+1$, (C1) holds but (C2) fails. For $l=3$, the
  children $\langle b,c,d,v,e\rangle$ and $\langle a,b,c,v,e\rangle$
  (deleting $a$ and $d$) share their facet at positions $2$ and $0$. That
  pair is not contained in the next split face $\{0,1,4\}$.
\end{itemize}

\section{The prefix family}\label{sec:prefix}

For $1\le k\le d-1$ and $l\in\{k,\dots,d\}$ put
\[
  F_l=\{0,1,\dots,k-1\}\cup\{l\},\qquad l(\sigma)=d-(\level(\sigma)\bmod(d-k+1)),
\]
split $\sigma[F_{l(\sigma)}]$, a $k$-face, and insert the new vertex at
$l(\sigma)$. The case $k=1$ is Maubach bisection, and $k=2$, $d=4$ is
\Cref{sec:face}. For $k=d-1$ there are two types, splitting the facets
$\partial_{d-1}$ and $\partial_d$ alternately.

\begin{theorem}\label{thm:prefix}
For every $d\ge2$ and $1\le k\le d-1$, the prefix scheme satisfies the
conclusions of \Cref{thm:face}, with invariants (A) and (B) for
$F(\sigma)=F_{l(\sigma)}$.
\end{theorem}

\begin{proof}
We check conditions (C0)--(C2) of \Cref{thm:cycle}. (C0) holds because
$l=\max F_l$. For (C1), the next split face is
$F_{l'}=\{0,\dots,k-1\}\cup\{l'\}$ with $l'=l-1$, or $l'=d$ if $l=k$.
Since $l\ge k$ and $l'\ne l$ (there are $d-k+1\ge2$ types), $l\notin F_{l'}$.
For (C2), let $p<q$ in $F_l$ and compute $\pos_p(q)$ and $\pos_q(p)$ as
in \eqref{eq:pos}.
\begin{itemize}
\item If $q<k$, the pair is $\{q-1,p\}$, contained in
  $\{0,\dots,k-2\}\subseteq F_{l'}$.
\item If $q=l$, the pair is $\{l-1,p\}$ with $p\le k-1$. When $l>k$,
  $l-1=l'\in F_{l'}$. When $l=k$, $l-1=k-1$ lies in the prefix.
\end{itemize}
This proves the theorem. It was also verified exhaustively for all
$d\le7$ and all $k$ by checking the two-cell lemma symbolically.
\end{proof}

\section{Typed cycles}\label{sec:cycles}

\begin{definition}
A \emph{typed cycle} in dimension $d$ is a finite sequence
$T=(\tau_0,\dots,\tau_{n-1})$ of pairs $\tau=(F_\tau,l_\tau)$, with
$F_\tau\subseteq\{0,\dots,d\}$, $|F_\tau|\ge2$, and $l_\tau\in\{0,\dots,d\}$.
A cell $\sigma$ has type $\tau(\sigma)=\tau_{\level(\sigma)\bmod n}$. It is
refined by the ordered stellar subdivision of $\sigma[F_{\tau(\sigma)}]$
with the new vertex at position $l_{\tau(\sigma)}$. We write
$F(\sigma)=F_{\tau(\sigma)}$ and $l(\sigma)=l_{\tau(\sigma)}$, and denote by
$\tau^+$ the type after $\tau$.
\end{definition}

A typed cycle may mix edge types ($|F|=2$), face types ($|F|=3$) and so on.
Maubach bisection is $((\{0,d\},d),(\{0,d-1\},d-1),\dots,(\{0,1\},1))$. The
invariants (A) and (B) of \Cref{def:face} make sense verbatim for any
typed cycle.

For $p\ne q$, let
\begin{equation}\label{eq:pos}
  \pos_p(q)=r+[\,r\ge l\,],\qquad r=q-[\,p<q\,],
\end{equation}
be the position of $\omega[q]$ in the child $c_p\omega$.

\begin{definition}\label{def:conditions}
A typed cycle is \emph{admissible} if for every type $\tau=(F,l)$:
\begin{enumerate}[label=(C\arabic*),start=0]
\item not $\min F<l<\max F$;
\item $l\notin F_{\tau^+}$;
\item for all $p<q$ in $F$, either $\pos_p(q)=\pos_q(p)$ or
  $\{\pos_p(q),\pos_q(p)\}\subseteq F_{\tau^+}$.
\end{enumerate}
\end{definition}

The conditions have the following meaning. (C0) makes two siblings meet
along the same \emph{ordered} facet. (C1) is the alternating property,
and it drives the direction lemma. (C2) says that two siblings satisfy (A).

\begin{lemma}[direction]\label{lem:dir}
Assume (C1) and the invariants, and let $\partial_i\sigma=\partial_j\omega$.
If $i\notin F(\sigma)$ then $\level(\omega)\le\level(\sigma)$, and if
$i\in F(\sigma)$ then $\level(\omega)\ge\level(\sigma)$. If the levels are
equal and $i\notin F(\sigma)$, then $i=j$ and
$\sigma[F(\sigma)]=\omega[F(\omega)]$.
\end{lemma}

\begin{proof}
If $\omega$ is finer, (B) for $(\omega,\sigma)$ gives $i\in F(\sigma)$. If
$\omega$ is coarser, (B) gives $i=l(\omega)$, and $l(\omega)\notin F(\sigma)$
by (C1), because $\tau(\sigma)=\tau(\omega)^+$. At equal levels with
$i\notin F$, (A) forces $i=j$. The two cells then differ only at
position $i\notin F$, so they have the same split face.
\end{proof}

\begin{lemma}[siblings]\label{lem:sib}
Let $\omega$ have type $\tau=(F,l)$ and let $p<q$ be in $F$. Under (C0), the
children $c_p\omega$ and $c_q\omega$ share the facet
$\partial_{\pos_p(q)}c_p\omega=\partial_{\pos_q(p)}c_q\omega$ as ordered
simplices. Under (C2), this pair satisfies (A).
\end{lemma}

\begin{proof}
Both children contain $v$ and $\omega\setminus\{\omega[p],\omega[q]\}$. The
vertices of $\omega$ keep their relative order, so the two ordered
facets can differ only in where $v$ sits. The vertex $v$ sits at position
$l$ in both children, so its predecessors are the first $l$ entries of
$\omega$ with $\omega[p]$ (respectively $\omega[q]$) removed.
\begin{itemize}
\item If $l\ge\max F$, these are $\{\omega[0..l]\}\setminus\{\omega[p]\}$ and
  $\{\omega[0..l]\}\setminus\{\omega[q]\}$.
\item If $l\le\min F$, both are $\{\omega[0..l-1]\}$.
\end{itemize}
In either case they agree after removing $\omega[q]$ and $\omega[p]$
respectively. If instead $p<l<q$, the vertex $\omega[l]$ precedes $v$ in
one facet and follows it in the other. The indices in the two children
are $\pos_p(q)$ and $\pos_q(p)$, and the children have type $\tau^+$, so (C2)
is exactly (A).
\end{proof}

\begin{lemma}[two cells]\label{lem:two}
Let $\sigma,\omega$ be adjacent with $\partial_i\sigma=\partial_j\omega$ and
$0\le\level(\sigma)-\level(\omega)\le1$, and let them satisfy the
invariants. Let $\tau=(F,l)$ be the type of $\omega$. Split $\omega$, and also
$\sigma$ if it contains $\omega[F]$. Then every adjacent pair formed by a
child of $\omega$ and $\sigma$ or a child of $\sigma$ satisfies the
invariants.
\end{lemma}

\begin{proof}
There are three cases.
\begin{enumerate}[label=(\alph*)]
\item \emph{Equal levels, $j\notin F$.} By \Cref{lem:dir}, $i=j$ and
  $\sigma$ has the same split face, so both cells are split. The children
  $c_p\sigma$ and $c_p\omega$ come from the same positional operation
  applied to two cells that differ only at position $i$, so they differ
  only at position $\pos_p(i)$. This satisfies (A) with equal indices.
  Children with different $p$ share at most $d-1$ vertices.
\item \emph{Equal levels, $j\in F$.} The facet omits $\omega[j]\in\omega[F]$,
  so $\sigma$ does not contain the split face and is not split. By
  \eqref{eq:outer}, $\partial_j\omega=\partial_l c_j\omega$. The pair
  $(c_j\omega,\sigma)$ has levels $\ell+1,\ell$. The finer index is
  $l=l(\sigma)$, and the coarser index is $i\in F(\sigma)$ by (A), since
  $j\in F$ and either $i=j$ or $\{i,j\}\subseteq F$. This is (B).
\item \emph{$\level(\sigma)=\level(\omega)+1$.} By (B), $i=l$ and $j\in F$, and
  only $\omega$ is split. Now $\partial_l\sigma=\partial_j\omega=\partial_l c_j\omega$,
  with equal levels and equal indices: (A).
\end{enumerate}
(OneLevel) holds in every case.
\end{proof}

It remains to see that there are no other new adjacencies. A facet of
$c_p\omega$ either omits $v$, in which case it is $\partial_l c_p\omega=\partial_p\omega$,
the outer facet treated in (b) and (c); or it contains $v$ and omits some
$\omega[s]$. If $s\in F$, the neighbor across it is the sibling $c_s\omega$
(\Cref{lem:sib}). If $s\notin F$, the facet contains all of $\omega[F]$
except $\omega[p]$, and the neighbor is the child $c_p\omega'$ of the cell
$\omega'$ across $\partial_s\omega$, which lies in the star of the split
face. This is case (a).

\begin{theorem}\label{thm:cycle}
Let $T$ be an admissible typed cycle and $K$ a conforming mesh satisfying
(OneLevel), (A) and (B), for example a strongly compatible seed. Then the
generalized Maubach algorithm (\Cref{alg:cycle}) terminates, only splits
cells of level at most $\level(\sigma)$, and preserves conformity and
the invariants. It is consistent, alternating and level-homogeneous.
\end{theorem}

\begin{proof}
The proof of the Maubach theorem of the article (\S\,Maubach's
scheme) carries over word for word, with \Cref{lem:dir,lem:sib,lem:two}
in place of its direction, children and two-cells lemmas. We only list
the replacements.
\begin{itemize}
\item The facets explored by the breadth-first search are those that
  contain the split face, $\partial_i c$ with $i\notin F(c)$.
\item The neighbor across such a facet is never finer
  (\Cref{lem:dir}).
\item A coarser neighbor $\omega$ meets $c$ along $\partial_j\omega$ with
  $j\in F(\omega)$, which passes whole to the child $c_j\omega$
  (case (c)).
\item Alternation is (C1).
\end{itemize}
\end{proof}

\begin{algorithm}[t]
\caption{$\textsc{CycleSubdivide}(K,\sigma)$ (\textsc{MaubachSubdivide} with $F$)}
\label{alg:cycle}
\begin{algorithmic}[1]
\State $\epsilon\gets\sigma[F(\sigma)]$;\ $\mathit{st}\gets[\,]$;\ $Q\gets(\sigma)$
\While{$Q\neq\varnothing$}
  \State $c\gets\mathrm{pop}(Q)$
  \If{$c$ not visited}
    \State mark $c$ visited; append $c$ to $\mathit{st}$
    \For{$i\in\{0,\dots,d\}\setminus F(c)$} \Comment{facets of $c$ containing $\epsilon$}
      \State $c_a\gets$ the cell across $\partial_i c$ (skip if none)
      \If{$\level(c_a)<\level(c)$}
        \State $\textsc{CycleSubdivide}(K,c_a)$;\ $c_a\gets$ the cell across $\partial_i c$
      \EndIf
      \State \textbf{if} $c_a$ not visited \textbf{then} push $c_a$ onto $Q$
    \EndFor
  \EndIf
\EndWhile
\State apply $(\epsilon,v)_{l(\sigma)}$ to every cell of $\mathit{st}$
\end{algorithmic}
\end{algorithm}

\paragraph{Necessity (verified, not proved).}
In dimension $4$ we ran an exhaustive symbolic check of the sibling and
two-cell configurations allowed by (A) and (B). Admissibility agrees with
this check on all $125$ one-type and $15625$ two-type cycles over
$|F|\in\{2,3,4\}$, and on $200000$ random cycles of length $3$--$5$. So,
for the invariants (A) and (B), conditions (C0)--(C2) are not only
sufficient but also necessary for these to be preserved. Other invariants
could in principle admit other schemes.

\section{Geometry}\label{sec:geometry}

\subsection{Faces alone never shrink cells}

\begin{proposition}\label{prop:edges}
If every type of a typed cycle has $|F|\ge3$, no edge of the mesh is ever
removed. In particular, the edges of the seed persist at every level, and
the maximal cell diameter does not tend to zero.
\end{proposition}

\begin{proof}
The stellar subdivision $(\delta,v)$ removes exactly the simplices
containing $\delta$. An edge does not contain a face $\delta$ with at
least three vertices.
\end{proof}

For the face scheme on Kuhn's cube with $v$ at the centroid of the face,
the child $\langle a,c,d,e,v\rangle$ of a level-$0$ cell keeps the diagonal
$ae$. The minimal quality $\mathrm{vol}/h^4$ falls by a factor $3$ per level
(verified): $0.33$, $0.11$, $0.037$, \dots, $0.0014$ at level $6$,
relative to level $0$. In the Riemann-surface application, the cells
around the curve become needles along seed edges and never get small.

\subsection{What to ask of a scheme on $\CP^2$}

On Kuhn's cube, Maubach bisection has finitely many similarity classes.
This is \emph{not} the criterion that matters on $\CP^2$: with
Fubini--Study midpoints no two cells are similar, and the classes are
infinite even for Maubach. What we want is:
\begin{enumerate}[label=(G\arabic*)]
\item \emph{Shrinking:} the maximal diameter $h_\ell$ at level $\ell$ tends
  to zero at a uniform rate. A necessary combinatorial condition is a
  finite \emph{edge lifetime}: there is an $L$ such that no edge of a
  cell survives in its descendants $L$ levels below.
\item \emph{Shape regularity:} the quality of the cells stays bounded
  below, or at least degrades slowly.
\end{enumerate}
The flat test on Kuhn's cube is still the right \emph{necessary} test. At
fine levels a cell of a smooth triangulation of $\CP^2$ lies in a region
where the Fubini--Study metric is flat up to $O(h^2)$, and geodesic
midpoints are affine midpoints up to $O(h^2)$. A scheme whose descendants
degenerate in the flat model therefore degenerates on $\CP^2$ as well.
Bounded quality on Kuhn's cube, and not finiteness of the classes, is the
property that transfers. We measure the \emph{shrinking exponent}
$\alpha$ defined by $h_\ell\approx N_\ell^{-\alpha}$, where $N_\ell$ is the
number of descendants of one cell. The ideal is $\alpha=1/d$: Maubach has
$\alpha=1/4$ exactly, since each cycle doubles the resolution with
$2^4$ cells.

\subsection{How a hybrid works}

A hybrid is an admissible typed cycle that contains both edge types and
face types. Refinement works as in Maubach's algorithm:
\begin{itemize}
\item Each cell knows its level, hence its type $\tau=(F,l)$.
\item To refine $\sigma$, collect the star of the split simplex
  $\sigma[F]$, recursively refining coarser neighbors first.
\item Apply one stellar subdivision with the new vertex at position $l$.
\end{itemize}
An edge step splits the star of an edge into $2$ children per cell. A
face step splits the star of a triangle, which in a $4$-manifold is a
cycle of cells since the link of a triangle is a circle, into $3$
children per cell. The new vertex goes at the midpoint of the edge or at
the centroid of the triangle. On $\CP^2$ we use the Fubini--Study
barycenter. The motivation for the Riemann-surface application is that
face steps split exactly the kind of simplex that is tested against the
curve.

The \emph{pointerless} (LPT) encoding of the article carries over. A cell
is still a root plus a path of child indices, now with $|F_\tau|$ possible
children at a step of type $\tau$. The neighbor-finding rules must be
rederived from \eqref{eq:outer} and \Cref{lem:sib}.

\subsection{Search results in dimension $4$}

All numbers below are from uniform refinement of one Kuhn simplex, with
midpoints and centroids. Descendants are deduplicated by similarity,
keeping the largest scale. This is exact while the number of classes is
below the cap, and a random sample above it. Sampling can only
\emph{overestimate} the minimal quality, so the degradation is at least
as bad as reported.
\begin{itemize}
\item \emph{Cycles of length $\le3$ mixing $|F|=2$ and $|F|=3$.} $2184$ are
  admissible ($48$ of length $2$, $2136$ of length $3$). Of these, $1392$
  have finite edge lifetime ($4$ to $7$ levels; Maubach has $3$ to $5$).
  None comes close to Maubach: after $12$ levels the best has
  $\alpha\approx0.11$ and minimal quality below $0.003$. This is expected,
  since a cycle of length $\le3$ cannot contain the $4$ edge types that
  Maubach needs.
\item \emph{Maubach with one face step.} Inserting one arbitrary face type
  into Maubach's $4$-cycle is never admissible. Replacing one edge type is
  admissible in exactly $3$ ways, all of which replace $(\{0,2\},2)$ by
  $(\{0,1,2\},l)$ with $l\in\{2,3,4\}$. The best is
  \[
    T^\ast=\bigl((\{0,4\},4),\ (\{0,3\},3),\ (\{0,1,2\},2),\ (\{0,1\},1)\bigr),
  \]
  a variable-prefix scheme ($k=1,1,2,1$). Per cycle of $24$ cells:
  \begin{center}
  \begin{tabular}{@{}lcccccc@{}}
  \toprule
  cycle & 1 & 2 & 3 & 4 & 5 & 6\\
  \midrule
  $h$, Maubach ($16\times$ cells/cycle)   & $0.500$ & $0.250$ & $0.125$ & $0.062$ & $0.031$ & $0.016$\\
  $q_{\min}$, Maubach                     & $1$ & $1$ & $1$ & $1$ & $1$ & $1$\\
  $h$, $T^\ast$ ($24\times$ cells/cycle)  & $0.500$ & $0.276$ & $0.176$ & $0.108$ & $0.068$ & $0.039$\\
  $q_{\min}$, $T^\ast$                    & $0.667$ & $0.298$ & $0.075$ & $0.022$ & $0.006$ & $0.002$\\
  \bottomrule
  \end{tabular}
  \end{center}
  Cells do shrink (finite edge lifetime, $\alpha\approx0.18$), but the
  minimal quality falls by a factor of about $3$--$4$ per cycle, so (G2)
  fails.
\end{itemize}

\subsection{Conclusions and open questions}

\begin{enumerate}
\item The face scheme is a correct stellar subdivision scheme, and so is
  the whole admissible typed-cycle family. The combinatorics of
  compatibility do not care whether the split simplex is an edge or a
  face.
\item Geometrically, pure face schemes are ruled out
  (\Cref{prop:edges}), and the hybrids we tested degenerate. Edge
  bisection is therefore needed for the \emph{geometry}, not for
  compatibility.
\item Open questions:
  \begin{itemize}
  \item longer hybrid cycles (length $\ge5$, several face steps per
    Maubach cycle);
  \item placing the face vertex somewhere other than the centroid, which
    cannot help pure face schemes (\Cref{prop:edges}) but might help
    hybrids;
  \item invariants other than (A) and (B).
  \end{itemize}
\item For the Riemann-surface application, the practical recommendation
  is to keep Maubach bisection for refinement and use the $2$-faces in
  the refinement \emph{criterion}, which is what the code already does.
\end{enumerate}

\section*{Scripts}
\texttt{scripts/} (Python 3, NumPy):
\begin{itemize}
\item \texttt{facesplit.py}: mesh simulator on the reflected Kuhn torus,
  with uniform and adaptive refinement.
\item \texttt{inv.py}: census of the $(\Delta\level,\type,i,j)$ pairs.
\item \texttt{symbolic.py}: two-cell check for the prefix family, $d\le7$.
\item \texttt{hybrid.py}: search for admissible edge+face cycles.
\item \texttt{charac\_rule.py}, \texttt{charac.py}: conditions (C0)--(C2)
  and their agreement with the exhaustive check.
\item \texttt{lifetime.py}: edge lifetime.
\item \texttt{geo.py}, \texttt{geo2.py}, \texttt{screen.py}, \texttt{insert.py},
  \texttt{traj.py}: geometry on Kuhn's simplex.
\end{itemize}

\end{document}
